% from ega4-2-fr.tex
\oldpage[IV-2]{200}

\item[\emph{(ii)}] Soit $B$ une $A$-algèbre essentiellement de type fini qui
soit un anneau local; si $\mathfrak p$ est l'idéal premier de $A$ au-dessus
duquel est l'idéal maximal de $B$, $B$ est aussi une
$A_{\mathfrak p}$-algèbre essentiellement de type fini (1.3.8); en vertu de
(i), on peut donc supposer que $\mathfrak p$ est l'idéal maximal $\mathfrak m$
de $A$. On a donc $B=C_{\mathfrak q}$, où $C$ est une $A$-algèbre de type
fini, $\mathfrak q$ un idéal premier de $C$ au-dessus de $\mathfrak m$; soit
$\mathfrak n\supset\mathfrak q$ un idéal maximal de $C$ (nécessairement
au-dessus de $\mathfrak m$); si l'on pose $k=A/\mathfrak m$,
$C/\mathfrak mC$ est une $k$-algèbre de type fini, donc le corps résiduel
$k'$ de $C$ en l'idéal maximal $\mathfrak n$ est fini sur $k$ (I, 6.4.2);
en vertu de (i), il suffira de prouver que $C_{\mathfrak n}$ est un
$\boldsymbol P$-anneau, puisque $C_{\mathfrak q}$ est un anneau local en un
idéal premier de $C_{\mathfrak n}$. On est ainsi ramené à démontrer le
lemme suivant.
\end{enumerate}
\end{proof}

\begin{lemma}[7.4.4.1]
\label{IV.2.s7.4.4.1-fr}
Soient $A$ un $\boldsymbol P$-anneau local noethérien, $k$ son corps
résiduel, $C$ une $A$-algèbre de type fini, $B$ un anneau local en un idéal
premier $\mathfrak n$ de $C$, tel que : 1\textsuperscript{o}
l'homomorphisme $A\to B$ soit local; 2\textsuperscript{o} le corps résiduel
$k'$ de $B$ soit une extension finie de $k$. Si $\boldsymbol P$ satisfait à
$(\mathrm P_{\mathrm{IV}})$, $B$ est un $\boldsymbol P$-anneau.
\end{lemma}

\begin{proof}
Soit $(x_i)_{1\leqslant i\leqslant m}$ un système de générateurs de la
$A$-algèbre $C$; montrons d'abord que l'on peut raisonner par récurrence sur
$m$. Soit $C'$ la sous-algèbre de $C$ engendrée par
$x_1,\ldots,x_{m-1}$, et soit $\mathfrak n'=\mathfrak n\cap C'$.
L'homomorphisme $A\to C_{\mathfrak n}$ se factorise en
$A\to C'_{\mathfrak n'}\to C_{\mathfrak n}$ et il est clair que
$A\to C'_{\mathfrak n'}$ et $C'_{\mathfrak n'}\to C_{\mathfrak n}$ sont
des homomorphismes locaux; si $k''$ est le corps résiduel de
$C'_{\mathfrak n'}$, $k\to k'$ se factorise de même en
$k\to k''\to k'$, donc $k''$ est extension finie de $k$ et $k'$ extension
finie de $k''$. L'hypothèse de récurrence entraîne que
$C'_{\mathfrak n'}$ est un $\boldsymbol P$-anneau; en outre, si
$S'=C'-\mathfrak n'$, $C_{\mathfrak n}$ est un anneau local de
$S'^{-1}C$; comme $C=C'[x_m]$, on a
$S'^{-1}C=C'_{\mathfrak n'}[x_m/1]$ et l'hypothèse de récurrence montre
encore que $C_{\mathfrak n}$ est un $\boldsymbol P$-anneau. On est ainsi
ramené au cas où $C$ est une $A$-algèbre engendrée par un \emph{seul}
élément $t$.

Appliquons (7.4.3, (ii)); en posant $A'=\widehat A$,
$B'=B\otimes_A A'$ est un anneau local de $C\otimes_A A'$ en un idéal
premier au-dessus de $\mathfrak n$; comme $A$ et $A'$ ont même corps
résiduel $k$, le corps résiduel de $C\otimes_A A'$ en cet idéal premier est
égal à $k'$. D'ailleurs $C\otimes_A A'$ est une $A'$-algèbre engendrée par
un seul élément. Il résulte donc de (7.4.3, (ii)) qu'on peut se borner à
démontrer (7.4.4.1) lorsque $\boldsymbol P$ est la propriété (iv$'$) de
(7.3.8), que $A$ est \emph{complet} et $C$ engendrée par un seul élément
$t$.

Pour montrer que les fibres formelles de $B=C_{\mathfrak n}$ sont alors
géométriquement régulières, appliquons le critère (7.3.16, b)). Soit $B_1$
une $B$-algèbre finie \emph{intègre}, donc engendrée par un nombre fini
d'éléments entiers sur $B$. En multipliant ces éléments par un élément de
$S=C-\mathfrak n$, on peut supposer qu'ils sont entiers sur $C$, et on peut
donc écrire $B_1=S^{-1}C_1$, où $C_1$ est une sous-$C$-algèbre de $B_1$
engendrée par un nombre fini d'éléments entiers sur $C$, donc une
$C$-algèbre \emph{finie et intègre}. D'autre part, $B_1$ est un anneau
semi-local, et tout anneau local $B_2$ de $B_1$ en un idéal maximal est un
anneau local de $C_1$ en un idéal premier, tel que $A\to B_2$ soit un
homomorphisme local; en outre, le corps résiduel de $B_2$ est une extension
finie de $k'$, donc de $k$. On voit donc (compte tenu de (7.3.14) et de (i))
qu'on est ramené à la question suivante : soit $C$ un anneau noethérien
\emph{intègre}, contenant un sous-anneau $C_0$ qui est une $A$-algèbre
engendrée par un seul élément $t$ et telle que $C$ soit une $C_0$-algèbre
\emph{finie}; si $\mathfrak n$ est un idéal maximal de $C$

